\section*{Methode und Grenzen}
\addcontentsline{toc}{section}{Methode und Grenzen}

\subsection*{Was gepr\"uft ist}

Jede Zahl aus einer Prim\"arquelle steht genau einmal in
\texttt{scripts/kennzahlen.py}, zusammen mit ihrer gedruckten Seite.
Abgeleitete Gr\"o{\ss}en werden gerechnet und nicht getippt, damit Text und
Tabellen einander nicht widersprechen k\"onnen. F\"unf Skripte erzwingen das:
\texttt{pruefe\_seiten.py} pr\"uft f\"ur jeden Rohwert, dass er auf der
zitierten Seite tats\"achlich vorkommt; \texttt{reihe.py} pr\"uft die
Zehnjahresreihe \"uber die Kette, dass die Vorjahresspalte eines Berichts dem
Berichtsjahr des Vorberichts gleicht; \texttt{check\_literals.py} verbietet
hartkodierte Zahlen im Flie{\ss}text; die beiden Fu{\ss}notenpr\"ufer halten
den Belegapparat konsistent.

Die Kettenpr\"ufung hat genau einen Bruch gefunden, und er war ein Befund und
kein Fehler: die Neuberechnung der St\"uckkosten des Jahres 2024
(Abschnitt~\ref{sec:margejeunze}). Auch die Seitenzerlegung tr\"agt einen
eigenen Wachhund: In der Annual Information Form 2017 erf\"ullte eine
Reservenzahl jedes Muster einer Seitenzahl. Sie wird verworfen, weil sie nicht
in die Seitenfolge des Dokuments passt -- eine Fu{\ss}zeilenform allein
gen\"ugt als Nachweis nicht.

\subsection*{Drei Farben, drei Verantwortlichkeiten}

Das Dokument trennt sichtbar, wer f\"ur eine Aussage einsteht.
\textbf{Schwarz} ist belegt: Zahl, Tabelle, Quelle, Seite. \bk{\textbf{Blau}
ist eine Einsch\"atzung im laufenden Text, die der Verfasser pr\"uft und
verantwortet.} \vd{\textbf{Rot} ist das Votum in Teil~\ref{sec:verdikt}: eine
aus den Befunden abgeleitete Empfehlung unter erkl\"arten Annahmen,
ausdr\"ucklich kein Beleg und keine Zusage.} Wer nur die schwarzen Stellen
liest, erh\"alt den gepr\"uften Sachstand ohne jede Wertung.

\subsection*{Prim\"ar und sekund\"ar}

Sechs Werte dieses Dokuments stammen nicht aus einem Prim\"ardokument,
sondern aus \"offentlich zug\"anglichen Zusammenstellungen: der
Analystenkonsens und die Kursreihen. Sie stehen in einem eigenen Block der
Zahlenschicht, werden vom Belegpr\"ufer nicht erfasst und sind in ihren
Tabellen als Sekund\"arquelle gekennzeichnet. Ratings und Kursziele stehen in
keinem Bericht eines Emittenten \"uber sich selbst; sie deshalb als
\enquote{nicht verf\"ugbar} zu f\"uhren w\"are eine unfertige Suche und keine
benannte L\"ucke gewesen.

\subsection*{Was nicht gepr\"uft ist}

Gepr\"uft ist die \emph{Herkunft} der Zahlen, nicht die Richtigkeit der
Schl\"usse. Vier Dinge leistet dieses Dokument ausdr\"ucklich nicht:

\begin{enumerate}
\item \textbf{Es prognostiziert den Goldpreis nicht.} Die Fortschreibung in
Teil~\ref{sec:flow} ist eine Kapazit\"atsrechnung unter erkl\"arten Annahmen.
Die entscheidende Gr\"o{\ss}e wird nicht gesch\"atzt, sondern in drei vom
Unternehmen selbst erzielten Auspr\"agungen durchgerechnet.
\item \textbf{Es gewichtet die Szenarien nicht.} Welches eintritt, ist eine
Einsch\"atzung \"uber den Goldmarkt. Die Prim\"arquellen tragen sie nicht, und
das Dokument nimmt sie nicht vorweg.
\item \textbf{Es bewertet die Lagerst\"atten nicht eigenst\"andig.} Reserven,
Gehalte und Kapitalwerte sind Angaben des Unternehmens und seiner
Sachverst\"andigen. Sie werden zitiert und nicht nachgerechnet; eine
lagerst\"attenkundliche Pr\"ufung liegt ausserhalb dieses Dokuments.
\item \textbf{Es ersetzt keine Anlageentscheidung.} Anlagehorizont,
Risikotragf\"ahigkeit und Portfoliozusammenhang stehen in keiner
Prim\"arquelle und kommen hier nicht vor. Das Votum in
Teil~\ref{sec:verdikt} nennt einen Schwellenkurs und beobachtbare Ausl\"oser;
ob und in welcher Gr\"o{\ss}e gehandelt wird, entscheidet der Verfasser.
\end{enumerate}

\subsection*{Wann diese Analyse veraltet}

Sie ruht auf dem Jahresbericht 2025 und dem Zwischenbericht zum
30.\,Juni 2026. Die tragende Gr\"o{\ss}e, der erzielte Goldpreis, \"andert
sich t\"aglich; die zweite tragende Gr\"o{\ss}e, Produktion und St\"uckkosten
gegen den ver\"offentlichten Prognosepfad, wird quartalsweise
ver\"offentlicht. Mit jedem neuen Zwischenbericht ist zuerst
Abschnitt~\ref{sec:margejeunze} fortzuschreiben, danach der Prognosevergleich
in Teil~\ref{sec:strategie} und erst dann zu pr\"ufen, ob das Votum noch
tr\"agt.
